%% chapters/10-wishart.tex
\chapter{The Wishart Distribution}
\label{ch:wishart}

\whereWeAre{%
When the data is Gaussian, the unscaled sample covariance matrix has a name: the \emph{Wishart distribution}. It is the multivariate generalization of the chi-square law, and it is the cleanest analytic object in this book.%
}

\PLACEHOLDER{Sources: existing main.tex Wishart sections (continuation Apr 16); Joseph thesis Section 3.2 Def 3.3.}

\section{Definition and density}
\PLACEHOLDER{$S = \sum X_i X_i^T$ for Gaussian $X_i$; the Wishart density formula.}

\section{Dimension one collapses to chi-square}
\PLACEHOLDER{Worked specialization. Why this is natural from the construction.}

\section{Dimension two: structure of the Wishart matrix}
\PLACEHOLDER{Diagonal $\chi_n^2$, off-diagonal $\sum x_i y_i$, basic moment computations.}

\section{Brief warm-ups}
\PLACEHOLDER{3--5 short questions.}

\chapterRecap{%
\begin{itemize}
  \item \PLACEHOLDER{Wishart is the multivariate chi-square.}
  \item \PLACEHOLDER{In dimension 1 it \emph{is} chi-square.}
  \item \PLACEHOLDER{In dimension 2 it has clean structure that powers Chapters 11--14.}
\end{itemize}%
}

\section*{Exercises}
\PLACEHOLDER{8--15 longer exercises.}
